\section{Síntese e Conclusão Geral}
\label{sec:conclusao}

Este documento perseguiu uma única pergunta ao longo de todas as suas
seções: entre os cinco arranjos institucionais que administram os
hospitais do SUS no estado de São Paulo (\textbf{Filantrópico, Público
Municipal, OSS, Administração Direta e Privado}), algum entrega
desempenho sistematicamente melhor nos cinco indicadores oficiais do
estudo? A resposta foi construída sobre o painel definitivo descrito na
Seção~\ref{sec:painel}: \textbf{314 hospitais acompanhados ano a ano de
2015 a 2025, 3.454 observações balanceadas}, destilado por um funil
auditável a partir de 830 estabelecimentos da base bruta do SIH/SUS. A
conclusão honesta é que a resposta depende radicalmente de qual pergunta
se faz, se quem opera mais intensamente, quem é mais eficiente ou quem
\emph{causa} melhor desfecho ao mudar de gestão, e que a evidência é
forte para umas e frágil para outras. Esta seção amarra as duas metades
do trabalho: a leitura descritiva do painel (Seções~\ref{sec:painel}
a~\ref{sec:qualidade}) e a leitura inferencial e de eficiência
(Seções~\ref{sec:estrategia} a~\ref{sec:fase2}).

\textbf{O que a fase descritiva estabeleceu.} A metade exploratória do
documento mapeou o terreno antes de qualquer modelo. Fixou os moldes
estatísticos de cada indicador (Seção~\ref{sec:distribuicoes}),
retratou o perfil bruto de cada categoria de gestão em medianas e
violinos (Seção~\ref{sec:categorias}) e mostrou que as redes diferem
tanto em desempenho quanto naquilo que as torna comparáveis, a
complexidade e o porte (Seção~\ref{sec:complexidadeporte}). Duas
ressalvas atravessam toda essa leitura e valem para o documento
inteiro. A primeira é que a categoria de gestão de cada hospital segue
a classificação assistencial do SIH adotada operacionalmente pela
equipe (a variável \texttt{modelo\_gestao\_proxy}), e \textbf{não um
cruzamento institucional depurado} caso a caso: é uma definição de
trabalho, não um cadastro jurídico. A segunda é que todo contraste
descritivo entre categorias reflete o \textbf{desempenho médio de quem
está em cada rede}, e não o efeito causal limpo do modelo de gestão,
pois hospitais não são sorteados para uma administração ou outra, e as
redes diferem de partida em porte, complexidade e vocação assistencial.
Os números descritivos situam; não julgam.

\textbf{O que a fase inferencial sustenta, por ordem de robustez.} A
metade causal e de eficiência do documento pode ser resumida não pela
ordem em que os métodos foram aplicados, mas pela \textbf{firmeza da
evidência que cada achado resistiu a produzir}. O que sobrevive a todos
os testes é o \textbf{regime assistencial mais intenso das OSS}: mais
produção, maior ocupação e maior uso de UTI (19,3\% acima da Direta,
p~igual a 0,003, na Seção~\ref{sec:inferenciafinal}), um contraste
transversal firme entre as redes, visível na produção física por leito
que não usa dinheiro nenhum (55 saídas por leito nas OSS contra 36 na
Direta) e replicado em todos os métodos, inclusive sem os conversores.
A esse traço soma-se a \textbf{maior eficiência produtiva das OSS}
(Seção~\ref{sec:fase2}), cerca de 25\% menos distância à fronteira que
a Direta, concordante entre DEA e SFA e igualmente estável à retirada
dos conversores: como a intensidade operacional, é uma característica
\emph{da rede}, medida na comparação entre hospitais distintos, e não um
efeito de conversão.

\textbf{O que é sugestivo, mas não definitivo}, é a queda de cerca de
\textbf{10\% no tempo médio de permanência} associada à gestão por OSS.
Diferentemente dos traços anteriores, esse é um efeito identificado
\emph{dentro} do hospital, pelos cinco conversores que trocaram a
Direta pela OSS (Seções~\ref{sec:efeitogestao}
e~\ref{sec:inferenciafinal}). Ele é consistente em direção e tamanho,
maior que o de 85\% dos conversores-placebo no teste de permutação e
resistente ao bootstrap selvagem por cluster, mas não alcança a
significância estatística convencional (p~igual a 0,153), porque
\textbf{cinco hospitais não bastam} para sustentar a afirmação com
rigor. É evidência sugestiva forte, não fato demonstrado. \textbf{O que
não se sustenta} é a queda de mortalidade após a conversão para OSS.
Visível nos dados brutos dos conversores, ela \textbf{não encontra
apoio no controle sintético}: em Sorocaba a lacuna frente ao gêmeo é
nula (0,54 ponto abaixo, p~igual a 0,67) e o próprio sinal inverte
conforme o painel usado no pool de doadores (era 1,09 acima antes da
ETAPA F) — fragilidade, não efeito; no Pérola Byington a queda existe mas
fica dentro do que sorteios de placebo produzem (p~igual a 0,32). A
Seção~\ref{sec:fase2} deu a esse ponto a demonstração mais direta:
retirando os cinco conversores da amostra, a diferença de mortalidade
entre as redes OSS e Direta simplesmente desaparece. A mortalidade
menor é um \textbf{fenômeno da trajetória dos hospitais que mudaram de
gestão, não uma diferença estrutural entre as redes}. \textbf{O que não
mostrou diferença em método nenhum} foi o faturamento real por saída, igual
entre as redes na leitura transversal (Seção~\ref{sec:resultados}) e na
hierárquica (Seção~\ref{sec:fase2}), e comum entre os placebos no teste
de permutação. \textit{Ressalva de leitura (regime de pagamento): hospitais OSS são custeados por contrato de gestão com orçamento global, não por reembolso de AIH --- o valor registrado por saída é, para eles, produção sem consequência financeira direta; a diferença de faturamento pode refletir o próprio regime de pagamento, a variável cujo efeito se busca medir, e não uso de recursos ou qualidade assistencial.} Sobre o grupo \textbf{Privado} (n~igual a 3), a regra do
documento foi mantida do início ao fim: entrou em todos os modelos
apenas como \textbf{controle de absorção}, para não contaminar as
demais categorias, e nunca recebeu interpretação substantiva, pois a
Seção~\ref{sec:fase2} mostrou seu efeito aparente encolher quase 90\%
assim que o modelo pondera a evidência pelo tamanho do grupo.

\textbf{Limitações do projeto como um todo.} Para além das ressalvas
técnicas discutidas seção a seção, quatro limitações operam no nível do
projeto e condicionam toda leitura dos resultados. A primeira, e a mais
importante, é que \textbf{toda a identificação causal repousa em apenas
cinco conversores}, três deles com um único ano de dados após a mudança
de gestão, uma base estreita para qualquer afirmação de efeito, por
mais cuidadosa que seja a metodologia. A segunda é a \textbf{ausência de
DRG no SIH}: a gravidade do caso, que idealmente ajustaria a comparação,
foi suprida por uma métrica própria de complexidade (de inspiração
Barcelona), cuja versão ponderada por mortalidade
(\texttt{complexidade\_pond\_mort}) permanece provisória. A terceira é
que o painel, por construção, retém apenas hospitais com mediana de
leitos superior a 50, o que \textbf{limita a generalização para
hospitais menores}, ausentes do estudo. A quarta é a \textbf{janela
temporal ainda curta} para os hospitais convertidos em 2025, cujo
pós-tratamento se resume a um único ano e não permite ver a trajetória
se estabilizar.

\textbf{Contribuição.} Com essas fronteiras claras, o que o projeto
entrega é, ao que se sabe, o \textbf{primeiro comparativo institucional
abrangente do estado de São Paulo}, e não apenas da Região
Metropolitana, a reunir, sobre o mesmo painel, três camadas
metodológicas que raramente convivem: a \textbf{estimação causal}
(efeitos aleatórios correlacionados de Mundlak, estudo de eventos,
o estimador de Callaway e Sant'Anna e controle sintético), a
\textbf{análise de eficiência} (DEA com correção de viés e fronteira
estocástica) e a \textbf{validação por modelos bayesianos
hierárquicos}. Mais do que o arsenal, a contribuição está no cuidado
metodológico explícito de \textbf{não confundir efeito de gestão com
efeito de trajetória de conversão}, distinção que, aplicada à
mortalidade, inverteu a leitura ingênua dos dados brutos e é
provavelmente o achado mais consequente do trabalho.

\textbf{Direções futuras.} Três frentes dão continuidade natural à
pesquisa. A primeira é a \textbf{fase qualitativa de entrevistas}, cuja
seleção de entrevistados e roteiro completo já constam da
Seção~\ref{sec:qualitativa}: as entrevistas semiestruturadas ali
desenhadas são dirigidas a investigar o mecanismo por trás do efeito no
tempo de permanência, pois, se ele é real, algo na gestão por OSS
encurta a internação, e entender o quê é o passo que os dados
administrativos não alcançam. A segunda é o acúmulo de \textbf{novos
anos de dados}, sobretudo para os três hospitais convertidos em 2025, o
único caminho para transformar a evidência dentro do hospital, hoje
sugestiva, em conclusão. A terceira é a \textbf{ratificação da forma
funcional da complexidade}, a escolha entre \texttt{complexidade\_pond\_mort}
e o fator de segurança \(s(M)\), ainda pendente de decisão junto à
pesquisadora principal, e que afeta o ajuste de gravidade em todas as
análises de mortalidade.

\textbf{Uma palavra final de cautela.} Os resultados aqui reunidos
foram construídos para serem lidos em conjunto e com seu contexto
metodológico completo. Isolado desse contexto, qualquer número deste
documento é vulnerável ao uso indevido, e dois pontos exigem alerta
expresso. Primeiro, \textbf{nada aqui autoriza a afirmação de que a
gestão por OSS reduz a mortalidade}: o controle sintético não encontra
descolamento significativo em nenhum dos dois conversores testados, a
comparação entre redes sem os conversores é nula, e citar a queda
bruta seria trair o próprio achado central do trabalho. Segundo, toda a
identificação causal apoia-se em \textbf{cinco hospitais convertidos},
uma base estatística frágil por natureza. Por isso, os achados deste
documento \textbf{não devem ser usados isoladamente por gestores
públicos para justificar decisões de modelo de gestão} sem o contexto
metodológico que os cerca. O que a pesquisa oferece com segurança é um
retrato: as OSS operam um regime mais intenso e mais eficiente; o que
ela expressamente não oferece é um veredito causal sobre qual modelo
salva mais vidas.
